This paper documents a repository-implemented pipeline that combines case retrieval, adaptive capability weighting, bottleneck detection, reflection, structured simulation evaluation, and standards-oriented exporters for simulated multi-domain operation planning. Using a single canonical evidence scope, the paper reports that the Full Memento--Hope--Reflection configuration outperforms the pure LLM baseline in the selected ablation study and in all five scenarios of the selected cross-scenario transfer suite.

Three additional contributions strengthen the original draft in response to reviewer feedback. First, a \texttt{BottleneckDetector} module replaces the static bottleneck-boost vector with a dynamic, history-aware diagnosis that identifies the current bottleneck stage and allocates capability emphasis proportionally to measured severity. The detector correctly flags \texttt{breach} as the persistent bottleneck across all iterations, and the weight-competition analysis reveals that breach and sustain share the \texttt{protection} capability dimension, creating a structural tension that a scalarized reward cannot fully resolve. Second, a memory-conditioned Hope coupling lets retrieved case metrics bias the fast-weight initialization, so the pipeline's two auxiliary modules share information through a common capability representation rather than operating as purely serial components. Third, the paper supplements the single-seed evidence with multi-seed replication across five seeds ($n=5$, paired $t(4)=5.593$, $p<0.01$, Cohen's $d=2.50$), confirming that the Full pipeline's gain over Pure LLM is statistically significant and robust under default HOPE parameters. A supplementary multi-seed experiment with optimized parameters ($\theta=0.9$, $\epsilon=0.005$, $n=3$) confirms the optimized Full outperforms Pure LLM in all tested seeds (mean $\Delta$MS = +0.0606), though seed 123 shows a reduced gain (+0.0262) relative to the default-parameter Full at the same seed (+0.0497), indicating that the optimal $(\theta,\epsilon)$ combination may interact with seed-specific initialization. External baseline comparisons (few-shot CoT, random search, single-pass Memento) bracket performance between a weak lower bound (MS=0.5189) and the optimized Full pipeline (MS=0.6707). Hyperparameter optimization of the HOPE SDE parameters ($\theta=0.9$, $\epsilon=0.005$) contributes a +0.0319 MS gain over the default-parameter Full setting, confirming that systematic parameter tuning is a cost-effective complement to architectural extension. A 20-iteration convergence analysis demonstrates that extended optimization yields further gains (+0.0869 MS over the 10-iteration result) and provides a mechanistic account of where improvement concentrates and where it stalls (\texttt{breach}). Cost--quality analysis uses recorded wall-clock timing data to inform practical deployment trade-offs.

Just as importantly, the paper makes its current boundaries explicit. It does not claim formal proofs, external benchmark leadership, or evidence beyond what the repository can support. The multi-seed evidence, baseline comparisons, and convergence analysis reported here substantially narrow the gap between the original source-traceable draft and the statistical rigor expected of a journal submission, but several limitations remain: (1) the cross-scenario study still uses a single seed (seed 7) and the reported confidence intervals reflect only within-seed Monte Carlo variance, not across-seed variability; (2) the Reflection-only cross-scenario negative result ($-0.0272$ mean MS deficit) has not yet been replicated across multiple seeds; (3) the per-component ablation is multi-seed only for Pure and Full settings; and (4) the case-bank size ablation (Section~4.10) is currently reported for a single seed. These limitations define the next refinement stage and do not invalidate the current contribution.
